\ficha{Obstfeld e Taylor (2003), Sovereign risk, credibility and the gold standard}{obstfeldtaylor2003}

\begin{identificacao}
OBSTFELD, Maurice; TAYLOR, Alan M. Sovereign risk, credibility and the gold
standard: 1870-1913 versus 1925-31. \textit{The Economic Journal}, v. 113,
n. 487, p. 241-275, 2003.

Artigo de história econômica quantitativa, 46 páginas, em inglês, lido na
íntegra. Versão lida: o working paper NBER 9345, de novembro de 2002, que já
traz na folha de rosto a nota de saída no \textit{The Economic Journal}. As
páginas citadas nesta ficha são as do working paper.
\end{identificacao}

\bloco{Problema e tese}

O câmbio fixo duro compra credibilidade? Os autores respondem comparando dois
períodos vizinhos de estabilização cambial, o padrão-ouro clássico e o
entreguerras, com a mesma medida de risco soberano. Antes de 1914 a adesão ao
ouro valia como sinal por si só e reduzia o custo de captação soberana em
Londres em cerca de 40 a 60 pontos-base. Entre 1925 e 1931 o sinal desaparece:
voltar ao ouro na paridade antiga não valia nada, quem voltou desvalorizado
ganhou algo, e o mercado passou a precificar dívida pública e pertencimento ao
Império Britânico.

\chave{shaved about 40 to 60 basis points from a country's external public borrowing cost}{p. 5}

\bloco{Argumento}

\begin{enumerate}[leftmargin=*,nosep]
\item O desenho corrige uma incomparabilidade. Bordo e Rockoff (1996) usaram
rendimentos de mercado secundário em Londres; Bordo, Edelstein e Rockoff (1999)
usaram emissões novas em Nova York, amostra pequena, com lacunas nos anos sem
emissão e viés de seleção, já que governo não lança título em condição
desfavorável (p. 2). Aqui a mesma medida se aplica aos dois períodos.

\item Antes de 1914 o resultado de Bordo e Rockoff se sustenta. Na amostra
original de sete países os autores reproduzem os mesmos 40 pontos-base (p. 16);
ampliada a amostra e acrescentados dívida sobre PIB, inflação, exportações e
renda per capita, o efeito fica em torno de 50 pontos-base e sobe a cerca de 65
na estimativa por painel dinâmico, sugerindo que a metodologia anterior
subestimava o prêmio (p. 24).

\item Dívida e inflação não são significantes antes de 1914, e esse é o achado
forte: o compromisso com o ouro bastava para o mercado presumir contenção
fiscal e monetária.

\chave{Gold was apparently a good enough seal of approval by itself}{p. 24}

\item O efeito Império não aparece antes de 1914: vale de 20 a 30 pontos-base
como estimativa pontual, menos que o ouro e sem significância (p. 19). Os casos
o desmontam pelos dois lados. A Austrália, membro do Império, paga prêmio de
cerca de 58 pontos-base por causa da crise de 1890; o Uruguai, fora do Império,
não paga. Pertencer ao Império não era condição necessária nem suficiente para
acesso preferencial a Londres (p. 20).

\item No entreguerras tudo se inverte. Retorno na paridade antiga rende zero;
retorno desvalorizado, caso da França, rende cerca de 35 pontos-base; dívida
passa a ser punida, com 10 pontos percentuais a mais de dívida sobre PIB
elevando o risco em 14,1 pontos-base; e o Império passa a valer entre 180 e 230
pontos-base (p. 30).

\chave{the gold standard strikes out but the Empire strikes back}{p. 30}

\item A leitura histórica é que o sinal do ouro era político, não técnico.
Depois de 1918, com partidos operários no jogo, a palavra do compromisso deixou
de bastar e os credores passaram a querer ver a contabilidade (p. 30 e p. 31).
\end{enumerate}

\bloco{Conceitos e definições operacionais}

Credibilidade é definida por um preço, não por um discurso: o desconto que
Londres concede ao título soberano, medido como spread sobre o consol. Para
isolar risco de calote e excluir risco cambial, só entram títulos longos, acima
de cinco anos, pagáveis em ouro ou em libras. Adesão ao ouro é variável binária
anual, majoritariamente de jure: exige conversibilidade legal e livre entrada e
saída de ouro, e quem impõe controle cambial conta como fora ainda que não
desvalorize (p. 35). Os autores admitem julgamento subjetivo na codificação,
sobretudo nos aderentes de fachada como Espanha e Itália (p. 14). Aparecem
ainda o pecado original e a regra contingente de Bordo e Kydland, que admitia
suspensão em guerra desde que restaurada a paridade anterior (p. 12).

\bloco{Evidência e método}

Painel anual de mais de 20 países, cotações de dezembro entre 1870 e 1913 e de
junho entre 1925 e 1931, até 871 observações antes da guerra e 143 no
entreguerras. Rendimentos do Global Financial Data, completados por
\textit{Investor's Monthly Manual}, \textit{The Economist}, \textit{Kimber's
Record} e \textit{Wall Street Journal}. Efeitos fixos de país com correção AR1 e
checagem por painel dinâmico de Arellano e Bond. Observações com spread acima de
mil a dois mil pontos-base, títulos já em default, são excluídas; sem isso o
efeito do ouro perde significância na amostra ampliada.

\bloco{Agentes e vertentes}

Bordo e Rockoff (1996) e Bordo, Edelstein e Rockoff (1999) são os
interlocutores diretos, o primeiro confirmado e o segundo contrariado.
Flandreau, Le Cacheux e Zumer (1998) acham efeito do ouro semelhante, de 35 a 55
pontos-base, mas encontram efeito forte da dívida já sob o padrão-ouro clássico
(p. 22). Do outro lado, a linhagem que privilegia o Império, de Marx e Hobson a
Davis e Huttenback, Edelstein e Ferguson, cuja hipótese de que a colônia trazia
garantia implícita de não calote, melhor selo que o ouro, é citada em bloco e
rejeitada para o período anterior a 1914 (p. 19). Eichengreen, Temin e Polanyi
sustentam a leitura política da ruptura de 1918.

\bloco{Críticas e limites}

Os autores pedem cautela extrema: a escolha do regime é endógena, e o
coeficiente não mede efeito puro do ouro, e sim benefício médio parcialmente
incondicional para países em condição de adotá-lo (p. 5). O painel entreguerras
tem no máximo sete observações por país, e o resultado anterior a 1914 depende
de truncar os spreads de calote. O texto discute Ferguson (2001, 2002); a versão
madura da tese do Império, em Ferguson e Schularick (2012), é posterior e reabre
o ponto aqui fechado, ao sustentar que para a periferia o ouro pouco valia e o
vínculo imperial valia muito. Sobre o Brasil o artigo quase nada diz: o país é
observação com efeito fixo alto, cerca de 195 pontos-base já descontado o ouro
(p. 20), com codificação vinda de Meissner (2002) e sem menção à Caixa de
Conversão. A única aparição do desenho institucional brasileiro está no
apêndice, na Caixa de Estabilização, cuja entrada no ouro é datada da lei de 18
de dezembro de 1926 e a saída de 7 de dezembro de 1929 (p. 35 e p. 36).

\begin{dialogo}
\item Procurar em 1906 o argumento de crédito externo em favor da Caixa, com o
vocabulário ``credito'', ``praça de Londres'', ``juro da divida''. Testa se o
mecanismo medido aqui, adesão ao ouro que barateia a captação soberana, era
categoria disponível ao debate ou reconstrução posterior da literatura.
\item Verificar se a defesa dos 15 pence contra o par legal de 27 pence invoca
sustentabilidade da paridade, e não só proteção ao café. O achado entreguerras
dá conteúdo teórico a essa linha, caso ela exista nas fontes.
\item Rastrear se dívida pública e déficit aparecem nos editoriais como
determinante do crédito externo. Antes de 1914 o mercado de Londres não punia
dívida; se a imprensa discute dívida com esse peso, mede-se a distância entre a
percepção local e o preço em Londres.
\item Na suspensão de agosto de 1914, verificar se os jornais a justificam como
suspensão legítima por guerra, com promessa de volta à mesma paridade. Seria
indício de que a regra contingente circulava como norma compartilhada.
\end{dialogo}

\usonocapitulo{Parte I. Fixa a magnitude e a natureza do prêmio de credibilidade
antes de 1914, cerca de 40 a 60 pontos-base, medido em custo de captação
soberana externa e não em estabilidade interna nem em volume de capital
atraído. Posiciona a dissertação entre Bordo e Rockoff (1996), aqui confirmado
com amostra maior e mais controles, e Ferguson e Schularick (2012), que retoma a
hipótese do Império rejeitada neste artigo. Serve também como advertência
metodológica, na voz dos próprios autores, de que o efeito estimado do ouro não
é causal. Não fornece evidência direta sobre a Caixa de Conversão, que o texto
não menciona.}
